% Windows Server Deployment — project report (generated documentation)
\documentclass[11pt,a4paper]{article}
\usepackage[margin=2.4cm]{geometry}
\usepackage[T1]{fontenc}
\usepackage{xcolor}
\usepackage{enumitem}
\usepackage{titlesec}
\usepackage[hidelinks]{hyperref}

\definecolor{noc}{HTML}{0F766E}
\definecolor{nocdark}{HTML}{134E4A}

\titleformat{\section}{\large\bfseries\color{nocdark}}{\thesection}{0.8em}{}
\titleformat{\subsection}{\normalsize\bfseries\color{noc}}{\thesubsection}{0.8em}{}

\newcommand{\ok}[1]{\textcolor{noc}{\textbf{#1}}}

\begin{document}

\begin{center}
  {\LARGE\bfseries Windows Server Deployment}\\[4pt]
  {\large Active Directory Forest, DNS/DHCP and Group Policy Baselines}\\[8pt]
  {\color{noc}\rule{\textwidth}{1pt}}\\[6pt]
  {\small Marouane Aabirrouche --- IT Infrastructure specialist --- R\'eseaux \& Syst\`emes}\\
  {\small Project report PRJ-02 \,\textbullet\, HomeLab documentation series}
\end{center}

\section*{Overview}
This project covers the design and deployment of a complete Windows Server
environment built around a full \emph{Active Directory} forest with integrated
\emph{DNS} and \emph{DHCP}, locked down by strict \emph{Group Policy Objects}
(GPOs) providing comprehensive user management and domain-wide security
baselines. The goal was to reproduce, end to end, the core identity and
configuration-management layer that a small-to-medium enterprise runs on,
and to validate that every control actually applies where it should.

\section{Objectives}
\begin{itemize}[leftmargin=1.6em]
  \item Deploy a full Active Directory forest from a clean Windows Server install.
  \item Integrate DNS and DHCP into the directory so addresses, names and
        domain membership are managed centrally.
  \item Enforce comprehensive user management through organizational units (OUs)
        and security groups.
  \item Apply strict GPOs forming a security baseline across the domain:
        password policy, account lockout, audit policy and hardening settings.
  \item Verify enforcement empirically (gpresult / RSOP) instead of assuming it.
\end{itemize}

\section{Environment}
The lab runs virtualized on the home server estate alongside the Docker/Linux
workloads described in the HomeLab report (PRJ-03). A dedicated Windows Server
virtual machine acts as domain controller; client machines join the domain to
exercise the policies.

\begin{table}[h]
\centering
\renewcommand{\arraystretch}{1.3}
\begin{tabular}{l l}
\textbf{Component} & \textbf{Role}\\
\hline
Windows Server VM & Domain controller: AD DS, DNS, DHCP\\
Domain & Dedicated lab forest, single domain tree\\
Clients & Domain-joined workstations for policy validation\\
Management & Server Manager, Group Policy Management Console (GPMC), RSOP\\
\end{tabular}
\end{table}

\section{Implementation}
\subsection{Forest and domain design}
The deployment starts from a clean installation: static addressing for the
server, then promotion to domain controller, which installs the Active Directory
Domain Services role and the integrated DNS zone. A single-domain forest keeps
the trust model simple while leaving room for future child domains. Default
containers were left untouched and all administration happens against a
deliberate OU hierarchy, so inheritance is predictable.

\subsection{Organizational units and user management}
OUs separate workstations, servers, users and service accounts. Security groups
map to functional roles rather than individuals, and permissions are granted to
groups only --- a pattern that keeps access reviews trivial. User accounts are
created centrally, with group membership driven by OU placement and role mapping.

\subsection{Integrated DNS and DHCP}
DNS lives inside the directory: the AD-integrated zone replicates with the
database itself, domain members resolve internal names reliably, and the server
forwards external queries upstream. DHCP is authorized in AD and serves the
lab scope with reservations where stable addressing matters; clients receive
the domain controller as their DNS server so domain joins and name resolution
``just work''.

\subsection{Group Policy baselines}
Policy is applied in layers --- domain-level baseline first, then OU-specific
refinements:
\begin{itemize}[leftmargin=1.6em]
  \item \textbf{Password and lockout policy}: length and complexity requirements,
        history, lockout thresholds and durations.
  \item \textbf{Account hardening}: restriction of local administrator use,
        controlled folder and service permissions.
  \item \textbf{Audit policy}: logon/account events and object access auditing
        feeding the event logs.
  \item \textbf{User experience restrictions}: drive mappings, desktop and start
        menu lockdowns where appropriate, demonstrating per-audience control.
\end{itemize}

\section{Validation}
Every GPO was validated with \texttt{gpresult} / Resultant Set of Policy on
target machines: password changes rejected when below baseline, lockout firing
after repeated failures, mapped drives appearing only for the intended groups,
and audit entries landing in the Security log. DNS resolution was tested for
internal records and forwarded lookups; DHCP leases, renewals and reservations
were confirmed from the client side.

\section{Outcomes and lessons learned}
\begin{itemize}[leftmargin=1.6em]
  \item Identity, name resolution and address assignment integrate cleanly when
        DNS/DHCP live with the directory --- split-brain issues disappear.
  \item GPO layering rewards discipline: baseline at the domain, exceptions at
        the OU, nothing enforced ad hoc on individual machines.
  \item Testing policy application (RSOP/gpresult) after each change catches
        scope mistakes immediately --- the difference between ``applied'' and
        ``should apply''.
\end{itemize}

\vfill
\begin{center}
  {\small\color{noc}\rule{0.5\textwidth}{0.6pt}}\\[4pt]
  {\footnotesize Generated documentation --- NOC portfolio, PRJ-02.}\\
  {\footnotesize Bouznika, Morocco \,\textbullet\, \today}
\end{center}

\end{document}
